\documentclass[11pt]{article}
\usepackage{icmlko}
\begin{document}

\title{유령 후보 --- 제일 커 보이던 짝은 그 도메인에 그 축이 없었다}
\author{월드모델 실험실}
\date{노트 503 · 2026년 8월 3일}
\maketitle

\begin{abstract}
노트 502 가 ``집 밖 짝이 안 갖는 축만 도메인별 가림이 안전하다''는 조항을
냈다. 그 조항이 지목한 안전한 사냥터(위키 $\cdot$ 태그 $\cdot$ gen
$\cdot$ grp)를 인구조사하고 제일 큰 후보들을 관문에 태웠다.

\textbf{두 가지가 한꺼번에 반증됐다 --- 하나는 바로 앞 노트의 조항이다.}

\begin{enumerate}
\item \textbf{\texttt{grp} 은 집 밖 짝에 $0\%$ 인데도 전이를 판다} ---
웹툰 \texttt{grp} 을 빼면 KR 만화 $-0.0106$($0/10$) · 비게임 앱
$-0.0100$($0/10$). \textbf{노트 502 조항은 하루도 못 갔다.}
\item \textbf{인구조사 후보 $22$ 짝 중 $10$ 짝은 그 도메인에 그 축이
아예 없다.} $1$ 등이던 팝업 \texttt{gen}($+0.0475$)은 팝업의
\texttt{gen} 덮음이 $\mathbf{0\%}$ 다. 실제로 팝업 \texttt{gen} 만 빼면
자 다섯이 \textbf{전부 정확히 $+0.0000$} --- 모형이 비트 단위로 같다.
\end{enumerate}
\end{abstract}

\begin{figure}[t]
\centering
\includegraphics[width=\textwidth]{figs/f1.pdf}
\caption{\textbf{가} --- 사냥터 크기. \textbf{나} --- 관문.
\textbf{다} --- 유령 후보.}
\end{figure}

\section{인구조사가 무엇을 재고 있었나}

인구조사는 축 하나를 \textbf{전 도메인에서} 빼고 도메인별 점수 차를
읽는다. 그러면 그 도메인에 그 축이 \emph{없어도} 수가 나온다 --- 다른
도메인에서 빠진 것이 공유 열을 바꾸기 때문이다(노트 498 조항: 한 칸을
바꾸면 모든 칸이 바뀐다).

$0\%$ 덮음인데 인구조사에 값이 실린 짝: 팝업 gen $+0.0475$ · 아이돌 grp
$+0.0094$ · 팝업 grp $-0.0152$ · 아이돌 gen $-0.0128$ · 도서 gen
$+0.0020$ · 시장팝업 gen $\cdot$ grp · 웹툰 gen · 펀딩 gen · 만화 grp.

\begin{quote}
\textbf{조항} --- \textbf{인구조사에서 후보를 뽑기 전에 그 도메인의 그 축
덮음을 본다.} $0\%$ 면 그 수는 \emph{이웃 효과}이고 후보가 아니다.
그 도메인에서 그 축을 빼는 일은 \textbf{정의상 아무 일도 안 한다.}
\end{quote}

이 조항은 노트 497 에도 소급된다 --- 거기 \texttt{target\_breadth} 표는
덮음이 다 $83 \sim 100\%$ 라 무사하지만, \emph{확인하고} 무사한 것이 아니라
\emph{운좋게} 무사했다.

\section{집 밖은 미리 못 짚는다}

노트 502 의 기계 설명은 그럴듯했다 --- 집 밖이 갖는 축을 줄이면 집 밖이
무너지고, 안 갖는 축을 줄이면 안 무너진다. 그런데 \texttt{grp}(집 밖
$0\%$)이 집 밖을 판다. 정리하면 이렇다.

\begin{center}
\begin{tabular}{lllr}
\hline
가림 & 집 밖 덮음 & 집 밖 결과 & \\
\hline
애니 트렌드 & $0\%$ & \textbf{산다} & KR $+0.0152$ \\
웹툰 \texttt{grp} & $0\%$ & \textbf{판다} & KR $-0.0106$ \\
세계애니 \texttt{grp} & $0\%$ & 비긴다 & KR $+0.0073$ · 앱 $-0.0017$ \\
웹툰 \texttt{tb} & $100\%$ & 판다 & KR $-0.0303$ \\
\hline
\end{tabular}
\end{center}

\textbf{덮음으로는 부호가 안 정해진다.} 집 밖 행은 그 축이 없어도 그 축의
\emph{갈래}를 지나가고(결측 쪽으로 떨어진다), 그 갈래 너머에 무엇이 있는지는
학습이 정한다.

\begin{quote}
\textbf{조항} --- \textbf{집 밖 효과는 미리 못 짚는다. 재야 안다.}
노트 502 의 사전 거르개는 폐기한다. 남는 것은 관문뿐이다.
\end{quote}

\section{그래도 하나 남았다}

\begin{center}
\begin{tabular}{lrrrrr}
\hline
팔 & 판 & 날짜 통제 & 무리 안 & KR & 앱 \\
\hline
$+$세계애니 \texttt{grp} & $+0.0018$ $8/10$ & $+0.0011$ $8/10$ & $+0.0011$ $6/10$ & $+0.0073$ $7/10$ & $-0.0017$ $4/10$ \\
\hline
\end{tabular}
\end{center}

자 다섯 중 넷이 양수고 앱만 $-0.0017$ 인데 부호가 $4/10$ 이라 방향이
없다. \textbf{통과도 기각도 아니다} --- 씨앗을 늘려 다시 보아야 하는
유일한 후보다.

\section*{남은 것}

\begin{itemize}
\item \textbf{세계애니 \texttt{grp} 을 씨앗 $20$ 으로} --- 앱 $-0.0017$
$4/10$ 이 잡음인지 아닌지가 전부다.
\item \textbf{덮음 $\geq 1\%$ 인 (도메인, 축) 짝만으로 인구조사를 다시
읽기} --- 이번 표의 절반이 유령이었으니 노트 497 $\cdot$ 503 의 후보
목록을 다시 세워야 한다.
\item \textbf{두 노트 연속으로 조항을 냈다가 다음 노트에서 죽였다}(노트
$493 \to 494$, $502 \to 503$). \textbf{기계 설명이 그럴듯할수록 빨리
쓰고 있다.} 조항을 낼 때 \emph{반증 예측}을 같이 적어야 한다.
\end{itemize}

\end{document}
